\subsection{Conventions for equivariant cohomology}\label{sec_pre_equ_coho}

Unless specified otherwise, all cohomology groups are taken with integral coefficients
in the sense of Borel equivariant cohomology.
Let $G = \mathrm{GL}_n(\C)$, let $T \subset G$ be the maximal torus of diagonal matrices,
and let $S = \C^\times$ act on the one-dimensional representation $\C_1 \cong \C$ by
the standard character of weight $1$ (loop rotation).
For a linear algebraic group $K$, we write $\Lambda_K = H_K^*(\pt; \Z)$.
Our primary coefficient rings are
\begin{equation}\label{eq:coefficient_rings}
 \Lambda_{T\times S} = \Z[t_1,\dots,t_n,\epsilon], \qquad
 \Lambda_{G\times S} = \Z[c_1,\dots,c_n,\epsilon],
\end{equation}
where $t_a\in H^2_T(\pt;\Z)$ is the equivariant first Chern class of the $a$-th coordinate line of $\C^n$ (corresponding to the character $\epsilon_a$), $\epsilon = c_1^S(\C_1)$, and $c_j = c_j^G(\C^n)$,
with gradings $\deg t_a = \deg\epsilon = 2$ and $\deg c_j = 2j$.
The restriction homomorphism $\Lambda_{G\times S}\hookrightarrow\Lambda_{T\times S}$ identifies
$c_j$ with the $j$-th elementary symmetric polynomial $e_j(t_1,\dots,t_n)$.
For an equivariant vector bundle $E$, we denote its equivariant Chern polynomial by
$c_u^K(E) = \sum_{j \ge 0} c_j^K(E) u^j$; this extends by multiplicativity to virtual bundles
as a formal power series in $u$.

\begin{lemma}\label{lem:borel}
Let $A$ be a complex algebraic torus (possibly trivial), and let $X$ be a
smooth complex quasi-projective variety equipped with an algebraic
$G\times A$-action.
Assume that the ordinary cohomology $H^*(X;\Z)$ is free over $\Z$ and
concentrated in even degrees.
Then the restriction homomorphism induces isomorphisms of graded algebras
\[
 H^*_{G\times A}(X;\Z)
 \otimes_{\Lambda_{G\times A}}\Lambda_{T\times A}
 \xrightarrow{\sim}H^*_{T\times A}(X;\Z),
 \qquad
 H^*_{G\times A}(X;\Z)
 \xrightarrow{\sim}H^*_{T\times A}(X;\Z)^W.
\]
Here the Weyl group $W = \mathfrak S_n$ acts on $H^*_{T\times A}(X;\Z)$ via the normalizer $N_G(T)$,
permuting the classes $t_a$, fixing $\Lambda_A$, and fixing all classes in the image of $H^*_{G\times A}(X;\Z)$.
\end{lemma}

\begin{proof}
Set $K = G \times A$ and $K_T = T \times A$.
Since $H^*(X;\Z)$ and $H^*(BK;\Z) \cong \Lambda_K$ are free $\Z$-modules
concentrated in even degrees, the Serre spectral sequence for the Borel fibration
$X \to EK \times_K X \to BK$ collapses at the $E_2$-page.
By the Leray--Hirsch theorem, any homogeneous integral basis $\{b_j\}$ of $H^*(X;\Z)$
lifts to a free $\Lambda_K$-basis $\{\widetilde b_j\}$ of $H_K^*(X;\Z)$.
These lifts restrict to a $\Lambda_{K_T}$-basis of $H_{K_T}^*(X;\Z)$,
yielding the base-change isomorphism.

Since the action of $W$ on $EK \times_{K_T} X$ covers the identity map on $EK \times_K X$,
each restricted basis element $\operatorname{res}(\widetilde b_j)$ is $W$-invariant.
Expanding an arbitrary class $\alpha \in H_{K_T}^*(X;\Z)$ in this basis, we see that $\alpha$
is $W$-invariant if and only if its coefficients lie in
$\Lambda_{K_T}^W = \operatorname{res}(\Lambda_K)$.
This identifies $H_K^*(X;\Z)$ with the invariant subalgebra $H_{K_T}^*(X;\Z)^W$.
\end{proof}

We adopt the convention that the projective bundle $\PP(E)$ parametrizes one-dimensional
linear subspaces (lines) in a vector bundle $E$ of rank $r$ over $X$.
Setting $x = -c_1^K(\CO_{\PP(E)}(-1))$, the equivariant projective bundle formula gives
\begin{equation}\label{eq:projective_bundle_formula}
 H_K^*(\PP(E);\Z) \cong
 H_K^*(X;\Z)[x]\Big/\Big(\sum_{j=0}^r c_j^K(E)x^{r-j}\Big).
\end{equation}
The classes $1, x, \dots, x^{r-1}$ form a free basis of $H_K^*(\PP(E);\Z)$ as an
$H_K^*(X;\Z)$-module, and the fiber integration map $\pi_*:H_K^*(\PP(E);\Z)\to H_K^{*-2(r-1)}(X;\Z)$
is normalized by $\pi_*(x^{r-1}) = 1$ and $\pi_*(x^i) = 0$ for $0\leq i < r-1$
(cf.\ \cite[Proposition~2.6.1]{anderson2024equivariant}).
